%RH LaTex Macros Apr 2014
%% query in the margin macro
\newcommand{\query}[1]{{$\bullet$}\marginpar{%
\vskip-\baselineskip %raise the marginpar a bit
\raggedright\footnotesize
%\itshape\hrule\smallskip#1\par\smallskip\hrule}}
\itshape\hrule\smallskip\parbox{.75in}{{$\bullet$} #1}\par\smallskip\hrule}}
\newcommand{\removequeries}{\renewcommand{\query}[1]{}}

%% web related requires: hyperref package
\newcommand{\webref}[2]{\href{#1}{\underline{#2}}}
\newcommand{\web}[1]{\webref{#1}{#1}}

%
% alphalist:  a), b), ...
%
\newcounter{alphanum}
\newenvironment{alphalist}{\begin{list}{\alph{alphanum})}
        {\usecounter{alphanum}\setlength{\leftmargin}{1.5em}}
        \rm}{\end{list}}

%% Grammar specific commands
\newcommand{\emptystring}{{\Large \tsym{\epsilon}}~}
\newcommand{\production}{\item}
\newcommand{\arr}{$\rightarrow$~}    %% arrow
\newcommand{\key}[1]{{\sffamily\bfseries #1}}  %% keyword
\newcommand{\tok}[1]{{\scshape\bfseries #1}}  %% token returned such as number
\newcommand{\non}[1]{{\it #1}}       %% nonterm
\newcommand{\sep}{~$\mid$~~}     %% alternation
\newcommand{\tsym}[1]{$\boldsymbol{#1}$}  %% token returned such as number

%% Walsh and Bit specific
%\newcommand{\contreq}[3]{#3 = {\cal C} (#1, #2)}
%\newcommand{\contr}[2]{{\cal C} (#1, #2)}
%\newcommand{\widevec}[1]{\overrightarrow{#1}}  too big and ugly
%\newcommand{\wt}{T}
\DeclareMathOperator*{\acos}{acos}
\DeclareMathOperator*{\asin}{asin}
\DeclareMathOperator*{\atan}{atan}
\DeclareMathOperator*{\argmax}{\arg\!\max}
\DeclareMathOperator*{\argmin}{\arg\!\min}
\DeclareMathOperator*{\bc}{bc}
\DeclareMathOperator*{\degray}{degray}
\DeclareMathOperator*{\deungray}{deungray}
\DeclareMathOperator*{\dimof}{dim}  % \dim is already defined
\DeclareMathOperator*{\extract}{extract}
\DeclareMathOperator*{\fitness}{fitness}
\DeclareMathOperator*{\gray}{gray}
\DeclareMathOperator*{\maxdim}{maxdim}
\DeclareMathOperator*{\pack}{pack}
\DeclareMathOperator*{\parity}{parity}
\DeclareMathOperator*{\pred}{pred}
\DeclareMathOperator*{\ungray}{ungray}
\DeclareMathOperator*{\unpack}{unpack}
\newcommand\degrees[1]{\ensuremath{#1^\circ}}
\newcommand{\betterthan}{\succ}
\newcommand{\binary}[1]{\mbox{\tt #1$_2$}}
\newcommand{\bits}[1]{\mbox{\tt #1}}
\newcommand{\bspace}{{\cal B}}
\newcommand{\chisq}{\mbox{$\chi^2$}}
\newcommand{\contains}{\supseteq}
\newcommand{\coversnot}{\not\supseteq}
\newcommand{\covers}{\supseteq}
\newcommand{\domain}{{\cal D}}
\newcommand{\dspace}{{\cal D}}
\newcommand{\embed}[2]{{\cal E} (#1, #2)}
\newcommand{\fhat}{\widehat{f}}
\newcommand{\forallsp}{\tnygap\forall~~}
\newcommand{\func}[1]{\mbox{\sl #1}}
\newcommand{\ghat}{\widehat{g}}
\newcommand{\given}{\,|\,}
\newcommand{\hhat}{\widehat{h}}
\newcommand{\isin}{\subseteq}
\newcommand{\isnotin}{\not\subseteq}
\newcommand{\logand}{\wedge}
\newcommand{\logiff}{\leftrightarrow}
\newcommand{\logimp}{\rightarrow}
\newcommand{\loglshift}{\stackrel{\mbox{\tiny shift}}{\longleftarrow}}
\newcommand{\logminus}{-}
\newcommand{\logeq}{=}
\newcommand{\lognand}{\uparrow}
\newcommand{\lognor}{\downarrow}
\newcommand{\logneg}{\sim\!\!}
\newcommand{\lognot}[1]{\overline{#1}}
\newcommand{\logor}{\vee}
\newcommand{\logpk}{\vdash}
\newcommand{\logrshift}{\stackrel{\mbox{\tiny shift}}{\longrightarrow}}
\newcommand{\logunpk}{\dashv}
\newcommand{\logxor}{\oplus}
\newcommand{\ones}{\vec{\tt 1}}
\newcommand{\pbits}[1]{{{\tt #1}'s}}
% \newcommand{\smallbc}{\mbox{\scriptsize\sl bc}}
% \newcommand{\smallpack}{\mbox{\scriptsize\sl pack}}
% \newcommand{\smallparity}{\mbox{\scriptsize\sl parity}}
\newcommand{\smallpm}{\raisebox{.3ex}{$\scriptstyle \pm$}}
% \newcommand{\smallunpack}{\mbox{\scriptsize\sl unpack}}
% \newcommand{\smbc}{\mbox{\scriptsize\sl bc}}
% \newcommand{\smpack}{\mbox{\scriptsize\sl pack}}
% \newcommand{\smparity}{\mbox{\scriptsize\sl parity}}
% \newcommand{\smpm}{\raisebox{.3ex}{$\scriptstyle \pm$}}
% \newcommand{\smunpack}{\mbox{\scriptsize\sl unpack}}
\newcommand{\strictcontains}{\supset}
\newcommand{\trm}[2]{\walsh_{#2}(#1)w_{#2}}
\newcommand{\unembed}[2]{{\cal C} (#1, #2)}
\newcommand{\vose}{\phi}
\newcommand{\walsh}{\psi}
\newcommand{\wcsum}{\wc_{sum}}
\newcommand{\wc}{{\cal K}}
\newcommand{\worsethan}{\prec}
\newcommand{\wt}{\mathcal{W}}
\newcommand{\yspace}{{\cal Y}}
\newcommand{\zeros}{\vec{\tt 0}}

\newcommand{\wfrac}{\frac{1}{2^L}}
\newcommand{\sumhp}[3]{\sum_{{#1}\elemof h_{#2, #3}}}
\newcommand{\sumbitlen}[1]{\sum_{{#1}=0}^{L-1}}
\newcommand{\sumlen}[1]{\sum_{{#1}=0}^{2^L-1}}
\newcommand{\sumsize}[1]{\sum_{{#1}=0}^{2^L-1}}  % don't use this anymore
\newcommand{\prodlen}[1]{\prod_{{#1}=0}^{L-1}}
\newcommand{\invpow}[1]{\inv{2^#1}}

%% Math
% \newcommand{\openm}{\vspace*{-.06in}\begin{equation}}
% \newcommand{\closem}{\vspace*{-.04in}\end{equation}}
% \newcommand{\openm}{$\displaystyle}
% \newcommand{\closem}{$}
\renewcommand{\emptyset}{\varnothing}
\newcommand{\onetoone}{$1\!\!-\!1$}
\newcommand{\reals}{\mathbb{R}}
\newcommand{\integers}{\mathbb{I}}
\newcommand{\posintegers}{\mathbb{Z}}
\newcommand{\modintegers}[1]{{{\posintegers}_{#1}}}

% logical operators in text form
\newcommand{\ortxt}{{\sc or}}
\newcommand{\andtxt}{{\sc and}}
\newcommand{\nottxt}{{\sc not}}
\newcommand{\xortxt}{{\sc exclusive-or}}
\newcommand{\Xortxt}{{\sc Exclusive-or}}
\newcommand{\truetxt}{{\sc true}}
\newcommand{\falsetxt}{{\sc false}}

\newcommand{\NP}{\mbox{NP}}
\newcommand{\infinity}{\infty}

\newcommand{\avg}[1]{\langle #1 \rangle}
\newcommand{\ceil}[1]{{\lceil #1 \rceil}}
% \newcommand{\comb}[2]{{\binomal{#1}{#2}}}
\newcommand{\comb}[2]{{#1\choose#2}}
\newcommand{\conv}{\circ}
\newcommand{\cross}{\!\times\!}
\newcommand{\dotprod}{\cdot}
\newcommand{\ds}{\displaystyle}
\newcommand{\elemof}{\in\,}
\newcommand{\notelemof}{\not\in\,}
\newcommand{\floor}[1]{{\lfloor #1 \rfloor}}
\newcommand{\intersect}{\cap}
\newcommand{\bigintersect}{\bigcap}
\newcommand{\inv}[1]{\frac{1}{#1}}
\newcommand{\mapto}{{\rightarrow}\,}
% \newcommand{\mod}{\bmod}
\newcommand{\order}[1]{{\cal{O}}(#1)}
\newcommand{\rounddown}[2]{{\lfloor #1 \rfloor \raisebox{-1.0ex}{\scriptsize $#2$}}}
\newcommand{\rspace}{{\cal R}}
\newcommand{\achspace}{{\cal H}}
\newcommand{\smtxt}[1]{{\scriptstyle\rm #1}}
\newcommand{\st}{\,:\,}
\newcommand{\sttxt}{~~~s.t.~~~}
\newcommand{\subsetof}{\subseteq}
\newcommand{\suchthat}{~:~}
\newcommand{\supermath}[1]{\raisebox{2ex}{$\scriptscriptstyle #1$}}
% \newcommand{\therefore}{.\!\!\cdot\!\!.}
\newcommand{\transpose}[1]{{#1}^{\scriptscriptstyle \sf T}}
\newcommand{\transposep}[1]{{(#1)}^{\scriptscriptstyle \sf T}}
\newcommand{\sminv}{\raisebox{.5ex}{\Large $-1$}}
\newcommand{\smtranspose}{\raisebox{.5ex}{\large\sf T}}
\newcommand{\txt}[1]{\mbox{~~#1~~}}
\newcommand{\tx}[1]{\mbox{#1}}
\newcommand{\ul}[1]{\underline{#1}}
\newcommand{\union}{\cup}
\newcommand{\bigunion}{\bigcup}

%% Sections
%Theorems
\newcounter{theoremnum}
\newcommand{\theorem}[1]{{\refstepcounter{theoremnum}\vspace*{.05in}\noindent\bfseries \underline{Theorem} \thetheoremnum\ } (#1) \index{#1 Theorem}\index{Theorems!#1}
\itshape
}

% theorem proving support
\newcommand{\case}[1]{\vspace*{.05in}\noindent{\bfseries\upshape CASE\ #1:}}
\newcommand{\numcase}[1]{\vspace*{.05in}\noindent{\bfseries\upshape CASE\ #1:}}
\newcommand{\numexample}[1]{\vspace*{.05in}\noindent{\bfseries\upshape EXAMPLE\ #1:}}
\newcommand{\derivation}{\vspace*{.05in}\noindent{\bfseries\upshape Derivation: }}
\newcommand{\theoremend}{\upshape}
% \newcommand{\proof}{\upshape \vspace*{-.05in}\noindent{\bfseries Proof: }}
\renewcommand{\qedsymbol}{\upshape \hspace*{0.9\columnwidth} $\Box$\hspace*{-1ex}\raisebox{.5ex}{$\Box$}}
\newcommand{\diamondqed}{\hspace*{0.9\columnwidth} $\Diamond$}
\newcommand{\blackqed}{$\blacksquare$}

%Lemmas
\newcounter{lemmanum}
\newcommand{\lemma}{{\refstepcounter{lemmanum}\vspace*{.05in}\noindent\bfseries \underline{Lemma} \thelemmanum\ }
\itshape}

% Conjectures
\newcounter{conjecturenum}
\newcommand{\conjecturenonum}{{\refstepcounter{conjecturenum}\vspace*{.05in}\noindent\bfseries \underline{Conjecture}\ }}
\newcommand{\conjecture}{{\refstepcounter{conjecturenum}\vspace*{.05in}\noindent\bfseries \underline{Conjecture} \theconjecturenum\ }
\itshape}

% Corollaries
\newcounter{corollarynum}[theoremnum]
\newcommand{\corollarynonum}{{\refstepcounter{corollarynum}\vspace*{.05in}\noindent\bfseries \underline{Corollary}\ }}
\newcommand{\corollary}{\refstepcounter{corollarynum}\vspace*{.05in}\noindent{\bfseries \underline{Corollary} \thetheoremnum \alph{corollarynum}\ }
\itshape}
\newcommand{\lastcorollarynum}{\thetheoremnum \alph{corollarynum}}

\newcommand{\namedcorollary}[1]{{\refstepcounter{corollarynum}\vspace*{.05in}\noindent\bfseries \underline{Corollary} \thetheoremnum \alph{corollarynum}\ } (#1) \index{#1 Corollary}\index{Corollarys!#1}}


% Observations
\newcommand{\observations}{\noindent{\bfseries Observations: }}

% Examples
\newenvironment{example}{\begin{quote}{\bf Example:~}}{\end{quote}}

%% General
\newcommand{\note}[1]{{\large\bfseries Note: #1}}
%\newcommand{\note}[1]{\fbox{\ #1\ }}
\newcommand{\abeam}[1]{\rule{#1}{0.0in}}
\newcommand{\citehere}{\fbox{\ add cite here\ }}
\newcommand{\defit}[1]{{\bfseries #1\/}\index{#1}}
\newcommand{\downstrut}[1]{\rule[-#1]{0pt}{#1}}
\newcommand{\upstrut}[1]{\rule{0.0in}{#1}}
\newcommand{\hr}{\noindent\rule{\columnwidth}{0.1\baselineskip}\\}
\newcommand{\indexit}[1]{{#1}\index{#1}}
\newcommand{\numnd}[1]{{$ #1^{nd} $}}
\newcommand{\numrd}[1]{{$ #1^{rd} $}}
\newcommand{\numst}[1]{{$ #1^{st} $}}
\newcommand{\numth}[1]{{$ #1^\textrm{th} $}}

\newcommand{\neggap}{\hspace*{-2ex}}
\newcommand{\exgap}{\hspace*{1ex}}
\newcommand{\tnygap}{\hspace*{.125in}}
\newcommand{\smgap}{\hspace*{.25in}}
\newcommand{\gap}{\hspace*{.5in}}
\newcommand{\lggap}{\hspace*{1.0in}}
\newcommand{\hugegap}{\hspace*{2.0in}}

\newcommand{\negvgap}{\vspace*{-2ex}}
\newcommand{\tnyvgap}{\vspace*{1ex}}
\newcommand{\smvgap}{\vspace*{2ex}}
\newcommand{\vgap}{\vspace*{4ex}}
\newcommand{\lgvgap}{\vspace*{8ex}}
\newcommand{\hugevgap}{\vspace*{16ex}}

%% Cheap Slide making
%\newcommand{\graphic}[2]{\fbox{\rule{0.0in}{#1}\large\rm {\bfseries #2 }}}
%\newcommand{\graphic}[2]{\framebox[0.99\textwidth]{\rule{0.0in}{#1}\large\rm {\bfseries #2 }}}
%\newcommand{\graphic}[2]{{\rule{0.0in}{#1}\large\rm {\bfseries #2 }}}
%\newcommand{\help}[1]{{\large\rm {\bfseries [} #1{\bfseries ]} }}



\newcommand{\graphic}[2]{\makebox[0.99\textwidth]{\rule{0.0in}{#1}\large\rm {\bfseries #2 }}}
\newcommand{\defhead}[1]{{\huge\sf \underline{#1}~}}
\newcommand{\newslide}{\clearpage\Huge\sf}
\newcommand{\slidehead}[1]{\clearpage{\Huge\sf \underline{#1}{\vspace*{.2in}}}}
\newcommand{\slideheadextra}[1]{{\vspace*{-.2in}}{\Huge\sf \underline{#1}}}
\newcommand{\slidesub}[1]{{\huge\sf #1{\vspace*{.1in}}}}
\newcommand{\say}[1]{{}}
\newcommand{\slidenote}[1]{}
\newcommand{\smbits}[1]{\mbox{\large\tt #1}}
\newcommand{\subhead}[1]{{\vspace*{.1in}\Huge\sf #1{\vspace*{.05in}}}}
\newcommand{\titlepagebox}[1]{\clearpage{\Huge\bfseries\sf
\vspace*{.38\textheight}
\setlength{\fboxrule}{2pt}\setlength{\fboxsep}{5pt}
\centerline{\fbox{\usefont{OT1}{lcmss}{m}{n}\fontsize{30pt}{36pt}\selectfont #1}}}}

% example of common parameters
% \textwidth=7.5in
% \textheight=9in
% \columnsep=.25in
% \oddsidemargin=0in
% \topmargin=-.5in
% \pagestyle{empty}
% \pagestyle{plain}
